\section{Comparative Methodological Positioning}
\label{sec:comparative}

Table~\ref{tab:comparison} separates mathematical methods, learning frameworks, and organizational descriptions. They operate at different levels and should not be ranked as if each were a complete alternative to the others.

\begin{table}[tbp]
\centering
\renewcommand{\baselinestretch}{1}\small
\renewcommand{\arraystretch}{1.25}
\caption{Scope of comparison. The final column identifies what a UVIF evaluation must establish, rather than assigning unsupported superiority scores.}
\label{tab:comparison}
\begin{tabularx}{\linewidth}{@{}>{\raggedright\arraybackslash}p{2.9cm}YY@{}}
\toprule
Approach & Relevant capability & Requirement for a UVIF comparison \\
\midrule
Constrained and multi-objective optimization & Explicit trade-offs, constraints, multiple criteria & Show value beyond relabelling objectives or adapting weights \\
Safe reinforcement learning & Sequential learning with risk and safety considerations & Match observations, constraints, horizons, and training resources \\
Active inference & Generative-model treatment of action and information acquisition & Compare fully specified resource-aware models \\
Resource-rational analysis & Benefits of cognition relative to computation costs & Include reserve carry-over and equivalent planning opportunities \\
Allostatic regulation & Anticipatory regulation of internal conditions & Test an independently measurable five-component mapping \\
Cognitive architectures & Integrated mechanisms for memory, learning, and action & Specify and measure the cost of the added regulatory layer \\
\uvif{} & Proposed explicit five-component decomposition & Establish identifiability, tractability, and incremental predictive value \\
\bottomrule
\end{tabularx}
\end{table}

\subsection{Optimization and learning as possible implementations}
Equation~\eqref{eq:toy_potential} demonstrates that one valid UVIF instance is also an optimization system. A multi-objective interpretation, a constrained policy search, or a learning-based controller may therefore implement the proposed organization. UVIF should not be distinguished solely by the use of several considerations or by adaptive weighting, both of which can be represented within established mathematical frameworks.

Reinforcement learning provides methods for learning from sequential interaction, while safe RL addresses additional risk and constraint requirements \citep{garcia2015saferl}. The question for UVIF is whether its explicit assignment of regulatory roles improves specification, transfer, or failure diagnosis under a fair comparison. A hard constraint can be a stronger protection mechanism than a soft force; describing it as external does not make it less effective.

Data-driven models likewise use explicit training objectives, even when the learned representations are not readily interpretable. Such models could supply prediction, perception, or value estimates inside a UVIF implementation. The regulatory decomposition does not by itself render their representations interpretable or improve generalization.

\subsection{The intended level of novelty}
The proposed contribution is a modelling discipline: name the regulatory demands, assign measurable state-dependent mechanisms to them, distinguish costs from constraints, and evaluate persistence together with task performance. It remains possible that a smaller decomposition, an established architecture, or a single well-specified objective will provide an equally good or better account.

Demonstrating novelty requires more than broad applicability. It requires a case in which the decomposition supports a useful measurement, intervention, prediction, or design decision that is difficult to obtain from a comparably explicit baseline. The current paper motivates that programme and gives a consistent template; it does not establish that the requirement has been met empirically.
